\chapter{Hardware e instrumentação}\label{cap:hardware}
\section{Placas ESP32 e restrições de pinos}
O Mestre usa um ESP32-S3-WROOM-1-N16R8, configurado para 16\,MiB de flash e 8\,MiB de PSRAM Octal. Os satélites usam ESP32-C3 SuperMini Plus de 4\,MiB. O Mestre tem dois núcleos Xtensa LX7, e cada satélite tem um núcleo RISC-V. Os projetos usam 160\,MHz. Clock do processador, capacidade da flash e memória externa são propriedades diferentes: aumentar a primeira não cria espaço para programa ou histórico.

\begin{table}[htbp]\centering\small
\caption{Placas e recursos usados pelos três programas}\label{tab:modelos-esp}
\begin{tabularx}{\textwidth}{@{}p{2.5cm}p{4.2cm}X@{}}\toprule
Unidade & Placa/módulo previsto & Recursos e conferência necessária\\\midrule
Mestre & ESP32-S3-WROOM-1-N16R8 & Dois núcleos a 160\,MHz; flash de 16\,MiB e PSRAM Octal de 8\,MiB. Conferir o módulo soldado e a entrada de alimentação da placa.\\
Equipe e Visitantes & ESP32-C3 SuperMini Plus & Um núcleo a 160\,MHz e flash de 4\,MiB em cada placa. Confirmar fabricante, pinagem e capacidade física da variante adquirida.\\\bottomrule
\end{tabularx}\end{table}

O repositório contém mais de uma revisão de circuito e não identifica de forma inequívoca a PCB montada. A tabela de GPIO abaixo descreve \textbf{a atribuição no firmware}, que deve ser comparada com conectores e trilhas da placa real. Antes de energizar, conferir polaridade e tensão dos trilhos, continuidade do GND, proteções e ausência de curto. Não aplicar 12 ou 24\,V diretamente aos GPIOs: UART, I\textsuperscript{2}C, entradas dos encoders e demais sinais lógicos do ESP trabalham em 3,3\,V.

\begin{table}[htbp]\centering\small
\caption{GPIOs que exigem reserva ou cuidado no ESP32-S3}\label{tab:pinos-especiais}
\begin{tabularx}{\textwidth}{@{}p{2.6cm}X@{}}\toprule
GPIO & Motivo e decisão de projeto\\\midrule
26--32 & Normalmente ligados à flash e à PSRAM pelo módulo; não destiná-los ao chicote.\\
33--37 no N16R8 & Ocupados pela interface da PSRAM Octal do Mestre; não destiná-los a periféricos externos. Em outra variante de placa, consultar seu esquema.\\
0, 3, 45 e 46 & Pinos de configuração de partida (\emph{strapping}); níveis externos indevidos durante reset podem alterar o modo de boot. Usar somente após analisar a placa e os resistores de polarização.\\
19 e 20 & Usados pelo USB nativo/Serial-JTAG; reaproveitá-los como GPIO pode impedir gravação ou diagnóstico por essa porta.\\\bottomrule
\end{tabularx}\end{table}

Essas restrições vêm da \href{https://docs.espressif.com/projects/esp-idf/en/stable/esp32s3/api-reference/peripherals/gpio.html}{documentação de GPIO do ESP32-S3}; pinos de \emph{strapping} e USB não são universalmente proibidos, mas devem preservar a função de partida e manutenção exigida na placa. O clock máximo e os dois núcleos constam da \href{https://www.espressif.com/sites/default/files/documentation/esp32-s3-wroom-1_wroom-1u_datasheet_en.pdf}{folha de dados do módulo}.

\section{Mapa de pinos do mestre}
\begin{table}[htbp]\centering\small
\caption{Interfaces do mestre definidas no firmware}\label{tab:pinos-mestre}
\begin{tabularx}{\textwidth}{@{}p{3.0cm}p{2.4cm}X@{}}\toprule
Interface & GPIO do Mestre & Ligação e função\\\midrule
Pedal Hall & 5 (ADC1) & Entrada calibrada para validar a faixa e calcular o pedido do pedal.\\
Seletor de marcha e ré & 16, 15, 4 & Entradas ativas em nível baixo para L, H e ré. A tensão do banco vem do INA228.\\
Encoder esquerdo & 17 (A), 18 (B) & Entradas de pulsos do lado esquerdo; cada canal recebe pull-up externo a 3,3\,V.\\
Encoder direito & 8 (A), 9 (B) & Entradas independentes do lado direito, com a mesma referência GND e pull-ups.\\
I\textsuperscript{2}C & 7 (SDA), 6 (SCL) & Barramento de MPU6050, OLED, INA228 e MCP4725, configurado a 400\,kHz. Confirmar endereços e níveis na montagem.\\
LEDs & 10, 11, 12 & Vermelho, verde e azul. O aviso amarelo combina vermelho e verde.\\
Buzzer passivo & 38 & Saída de tons por PWM; dimensionar o estágio de acionamento para o buzzer instalado.\\
UART Visitantes & 19 (TX), 20 (RX) & Canal dedicado ao satélite público; cruzar TX com RX do outro ESP.\\
UART Equipe & 39 (TX), 40 (RX) & Canal dedicado ao satélite técnico; não unir as duas UARTs.\\
microSD SPI2 & 2 (CLK), 42 (MOSI), 1 (MISO), 41 (CS) & Barramento do cartão. Conferir alimentação, nível lógico e frequência efetiva reportada na montagem.\\
Modo A/B, ré e relé & 14, 13, 47, 21 & Saídas ativas em alto. O relé usa o contato normalmente aberto.\\\bottomrule
\end{tabularx}\end{table}

\section{Encoders e medição do movimento}
O projeto prevê um encoder incremental E38S6G5-600B-G24N em cada lado do veículo. Cada unidade fornece dois sinais de pulsos, A e B, defasados no tempo. O modelo documentado tem 600 pulsos por volta do próprio eixo. A defasagem permitiria deduzir o sentido de rotação com decodificação em quadratura; nesta revisão, o firmware conta as bordas dos dois canais para melhorar a resolução, mas \textbf{não calcula o sentido}. Por isso, recuar a roda também aumenta o percurso acumulado.

As saídas são NPN de \emph{coletor aberto}. O transistor interno do encoder apenas liga a saída ao GND quando conduz; desligado, ele deixa a saída eletricamente aberta. Sem um resistor de \emph{pull-up}, esse estado aberto ficaria indefinido. Cada canal A e B recebe, portanto, um resistor externo de 4,7\,k$\Omega$ até 3,3\,V: quando o transistor está aberto, o ESP lê alto; quando conduz, ele drena a pequena corrente do resistor e o ESP lê baixo. A alimentação do corpo do encoder é separada da tensão dos sinais. A documentação do modelo recomenda alimentação externa de 12\,V, mas o conector da revisão de esquema analisada sugere outra linha; conferir etiqueta, folha de dados e montagem antes de ligar. \textbf{Nunca ligar o pull-up dos sinais à alimentação de 12 V.} O GND dessa fonte e o GND do Mestre precisam ser comuns para que alto e baixo tenham referência definida.

O uso de \emph{dois encoders}, um em cada lado, permite estimar a distância média mesmo quando as rodas não percorrem exatamente o mesmo caminho em uma curva. Também permite comparar os lados: diferença persistente pode indicar deslizamento, geometria incorreta, interrupção de um cabo ou falha de um sensor. Essa comparação é diagnóstica e não identifica sozinha a causa. Cabos de sinal devem ser afastados dos de potência; confirmar blindagem, retorno de GND, forma de onda e ausência de contagem com rodas paradas.

O periférico PCNT conta as quatro bordas dos canais A e B de cada lado. Assim, 600 P/R produzem 2\,400 contagens por volta do eixo do encoder. O filtro de pulsos curtos configurado é de 1\,000\,ns. O odômetro é distância acumulada sem sinal, não posição orientada; tampouco decide a marcha do veículo.

Para cada lado $j$, a distância estimada é
\begin{equation}
d_j=\frac{n_j}{4P_j\,i_j}\,\pi D_j ,
\qquad
v_j=\frac{\Delta d_j}{\Delta t}\,3{,}6 ,
\label{eq:encoder}
\end{equation}
em que $n_j$ é a contagem acumulada, $P_j$ os pulsos nominais por volta, $i_j$ as voltas do encoder por volta da roda, $D_j$ o diâmetro efetivo em metros e $\Delta t$ o intervalo em segundos. O fator 3,6 converte m/s em km/h. O firmware configura $D=0{,}300$\,m e $i=1{,}0$ para cada lado; são valores a medir na montagem. O serviço usa a média das distâncias esquerda e direita, rejeita saltos acima do limite diagnóstico de 200\,km/h e alerta sobre divergência persistente. Esse limite é filtro de plausibilidade, \textbf{não} velocidade certificada do carro.

O PCNT acumula pulsos por hardware e o odômetro lê as duas rodas a cada 20\,ms. A velocidade divulgada usa a média. Para limitar tração, o Mestre usa a roda mais rápida. O percurso da sessão reinicia no boot, enquanto o total é recuperado da NVS. A cada 100\,m acumulados, o odômetro sinaliza a necessidade de persistir. Outra tarefa executa a escrita. Uma interrupção de energia pode perder o trecho que ainda não foi salvo. Para o ensaio, medir circunferência efetiva com o pneu carregado, confirmar a relação de transmissão e comparar ambos os lados com uma referência independente.

\section{MPU6050: aceleração e rotação}
O MPU6050 reúne um \emph{acelerômetro}, que responde à aceleração e à gravidade, e um \emph{giroscópio}, que mede a velocidade de rotação. Ele conversa com o Mestre por I\textsuperscript{2}C no endereço \texttt{0x68}; SDA vai ao GPIO 7 e SCL ao GPIO 6. A placa do sensor deve ficar rigidamente fixada ao chassi. O firmware foi parametrizado para X vertical positivo para cima e Y longitudinal positivo para frente. Se a montagem girar ou inverter o módulo, é preciso ajustar eixo e sinal em \path{firmware/mestre-s3/include/nucleo/configuracao/veiculo.h} e repetir a calibração.

O acelerômetro não distingue sozinho a aceleração do veículo da componente da gravidade produzida por inclinação da placa. O serviço amostra a 10\,Hz, com filtro interno de 5\,Hz, usa o giroscópio para acompanhar mudanças rápidas de orientação e corrige lentamente a estimativa de gravidade quando o movimento permite. Subtrai então a gravidade projetada no eixo longitudinal, filtra ruído e usa dois limiares com confirmação temporal para classificar aceleração, desaceleração ou condição próxima do repouso sem alternar a cada vibração. O valor em \emph{g} é convertido nominalmente por $a_{\mathrm{m/s^2}}=9{,}80665\,a_g$.

Ao iniciar, o veículo deve permanecer parado: o firmware usa 10 amostras para estimar os desvios e rejeita calibração com giro, intensidade de gravidade implausível ou eixo longitudinal quase vertical. Uma calibração aceita não corrige uma fixação que depois se moveu. A velocidade principal vem dos encoders, pois integrar a aceleração inercial para obter velocidade acumularia erro de inclinação, ruído e viés. Telemetria e diagnóstico distinguem sensor presente, calibração concluída e dado válido; conferir esses estados antes de comparar aceleração com outro instrumento.

\section{Tela, sinalização e armazenamento}
A OLED SSD1306 de 0,96\,pol e 128$\times$64 pixels usa SDA/SCL do mesmo barramento I\textsuperscript{2}C que atende o MPU6050. O firmware procura os endereços \texttt{0x3C} e \texttt{0x3D}; isso permite encontrar variantes comuns sem alterar o programa. Conferir alimentação e nível dos pull-ups no módulo adquirido: o barramento precisa permanecer compatível com 3,3\,V, e resistores de módulos em paralelo alteram sua resistência equivalente. A seção~\ref{sec:eletronica} explica o dimensionamento. A tela apresenta etapas de partida, estado, medições disponíveis e progresso de operações; sua ausência é reportada, mas não impede as UARTs nem o registro. As páginas e os avisos são detalhados na seção~\ref{sec:diagnostico}.

Os LEDs externos mostram verde para condição normal, amarelo para atenção ou instalação e vermelho para falha crítica. O GPIO 12 aciona o LED amarelo pelo sinal \texttt{LED\_BLUE\_CMD} do esquema. O LED RGB integrado, no GPIO 48, também distingue partida, Wi-Fi e as fases de OTA. O painel Desenvolvedor permite desligar o RGB integrado e salva a preferência no Mestre. Cada LED externo precisa de resistor em série dimensionado para o componente instalado. O buzzer passivo recebe tons por PWM no GPIO 38. Como o GPIO tem capacidade de corrente limitada, prever transistor e resistor de base conforme o buzzer escolhido, com retorno em GND. Cargas indutivas precisam de proteção contra o pulso de desligamento. A porcentagem de PWM configura a razão do pulso.

O microSD usa SPI2, com CLK, MOSI, MISO e CS nos GPIOs 38, 39, 40 e 41. Ele armazena CSV, logs e cópia de ajustes; uma tarefa de gravação recebe amostras por fila para que o tempo de escrita do cartão não dite o ritmo da aquisição. A falta do cartão elimina esses arquivos persistentes enquanto durar a falha, embora o fluxo ao vivo possa continuar. Confirmar alimentação, nível lógico, desacoplamento e pinagem no módulo comprado; um módulo com regulador ou conversores de nível pode ter requisitos diferentes de um encaixe de cartão simples. O capítulo~\ref{cap:memorias_google} descreve arquivos e recuperação.

\section{UART e referência elétrica entre os três ESPs}
\begin{table}[htbp]\centering\small
\caption{Ligações bidirecionais entre o Mestre e os satélites}\label{tab:uart-fios}
\begin{tabularx}{\textwidth}{@{}p{2.2cm}X X p{1.8cm}@{}}\toprule
Enlace & Mestre TX $\to$ satélite RX & Satélite TX $\to$ mestre RX & Referência\\\midrule
Equipe & GPIO 39 $\to$ GPIO 5 & GPIO 6 $\to$ GPIO 40 & GND comum\\
Visitantes & GPIO 19 $\to$ GPIO 5 & GPIO 6 $\to$ GPIO 20 & GND comum\\\bottomrule
\end{tabularx}\end{table}
Cada satélite tem um par TX/RX exclusivo do Mestre. TX transmite e deve chegar ao RX da outra placa; unir TX com TX impede a comunicação. Os três GNDs precisam compartilhar uma referência elétrica para que cada receptor interprete corretamente os níveis dos sinais. \textbf{Não é necessário interligar as linhas de 3,3\,V das três placas}: cada uma pode receber alimentação pelo seu ramo previsto, desde que os sinais e retornos respeitem os limites elétricos. Aterramento comum não significa usar o fio fino da UART como caminho de retorno da corrente de potência.

Ambos os enlaces usam 1\,000\,000 baud, oito bits de dados, sem paridade e um bit de parada (8N1). O Mestre envia telemetria e recebe confirmações; a Equipe também encaminha solicitações, e os satélites podem receber imagens OTA pelo mesmo enlace. Antes de interpretar erro de protocolo, conferir alimentação de cada placa, continuidade do GND, cruzamento TX/RX, nível alto de 3,3\,V e integridade do cabo. Se as placas estiverem fisicamente afastadas ou em ambiente ruidoso, verificar a forma de onda e o aterramento na montagem real.

\section{Eletrônica da placa e instrumentação elétrica}\label{sec:eletronica}
\subsection{Alcance do esquema analisado}
As folhas hierárquicas consultadas se identificam como revisão v4 de 18/10/2026. Elas representam o circuito projetado; o repositório não associa essa revisão a uma PCB comprovadamente montada no veículo. O símbolo do mestre usa a designação ESP32-S3-DevKitC, enquanto o firmware é configurado para o módulo N16R8. Antes de fabricar ou modificar a placa, conferir módulo, footprint, pinos físicos, esquema, lista de materiais e identificação impressa da PCB. Valores e ligações citados abaixo descrevem o esquema, salvo quando o texto identifica uma implementação no firmware ou um requisito de ensaio.

\subsection{Cadeia de alimentação: 48 V, 5 V e 3,3 V}
A folha de alimentação representa uma entrada rotulada 48 V, um módulo buck identificado como 100 V/5 V--2 A, saída de 5 V com C1 de 47\,$\mu$F e um estágio indicado como AMS1117 para produzir 3,3 V. O buck é um conversor chaveado abaixador: transfere energia por uma chave e um indutor para obter 5 V com perdas menores que uma queda linear direta desde o barramento do veículo. \textbf{Na distribuição projetada, a linha de 5 V alimenta os módulos ESP e outros consumidores de 5 V; a linha de 3,3 V fornece alimentação a determinados periféricos e os pull-ups dos sinais.} Os pinos GPIO dos ESPs continuam sendo sinais lógicos de 3,3 V, mesmo quando a placa de desenvolvimento recebe 5 V no pino de alimentação.

O rótulo de 100 V/5 V--2 A identifica o símbolo customizado, não uma medição da capacidade do módulo comprado. Confirmar fabricante, faixa de entrada, tensão máxima do conjunto de baterias carregado, corrente contínua e de pico, dissipação, proteção contra surtos e polaridade. Para o estágio linear, a potência térmica aproximada é $P_{\mathrm{LDO}}=(V_{\mathrm{in}}-V_{\mathrm{out}})I$; de 5 V para 3,3 V, cada 0,5 A requer dissipar cerca de 0,85 W. É preciso somar os consumos de todos os ramos, inclusive os picos de Wi-Fi e escrita no SD, e medir as duas linhas durante a partida. O fato de uma placa ESP aceitar 5 V na entrada não permite alimentar seus GPIOs com 5 V.

\begin{figure}[htbp]\centering
\begin{tikzpicture}[>=Latex,bloco/.style={draw=nexus,rounded corners,fill=pale,align=center,minimum height=10mm,font=\small}]
\node[bloco,text width=2.2cm] (bat) at (0,0) {Entrada rotulada\\48 V};
\node[bloco,text width=2.2cm] (buck) at (3.3,0) {Buck\\5 V};
\node[bloco,text width=3.0cm] (cinco) at (7.3,0) {Placas ESP e módulos\\alimentados em 5 V};
\node[bloco,text width=2.5cm] (ldo) at (3.3,-1.9) {Estágio AMS1117\\3,3 V};
\node[bloco,text width=3.0cm] (tres) at (7.3,-1.9) {Módulos de 3,3 V\\e pull-ups};
\draw[->,thick] (bat)--(buck);
\draw[->,thick] (buck)--(cinco);
\draw[->,thick] (buck)--(ldo);
\draw[->,thick] (ldo)--(tres);
\end{tikzpicture}
\caption{Distribuição de alimentação representada no esquema; confirmar cada saída na placa montada.}\label{fig:alimentacao}
\end{figure}

\begin{table}[htbp]\centering\small
\caption{Linhas de alimentação e cargas indicadas no esquema}\label{tab:alimentacao-modulos}
\begin{tabularx}{\textwidth}{@{}p{3.6cm}p{2.0cm}X@{}}\toprule
Conjunto & Linha indicada & O que conferir na montagem\\\midrule
Placas ESP mestre, Equipe e Visitantes & 5 V na placa & Entrada de 5 V e regulador interno da placa usada; GPIO e UART permanecem em 3,3 V.\\
Buzzer, circuito do relé e conector do módulo INA228 & 5 V & Corrente da bobina, módulo realmente instalado e compatibilidade dos sinais com 3,3 V.\\
MPU6050, OLED e microSD & 3,3 V & Tensões nos conectores, consumo e desacoplamento local.\\
SDA/SCL e canais A/B dos encoders & 3,3 V via resistores & Pull-ups são referências de nível lógico, não fontes de alimentação dos módulos.\\
Corpo dos encoders & Divergente & O conector do esquema indica 5 V; a documentação do modelo cita alimentação externa de 12 V. Conferir etiqueta e folha de dados do componente instalado antes de ligar.\\\bottomrule
\end{tabularx}\end{table}

O esquema apresentado não fecha, por si só, uma análise de proteção da entrada: fusível, proteção contra inversão, sobretensão transitória, desconexão de manutenção e segregação de cabos de tração precisam ser identificados na placa real e nos documentos do veículo. O GND comum é necessário para sinais UART e I\textsuperscript{2}C, mas correntes altas do motor não devem circular pelo retorno fino dos sensores. Planejar retornos e pontos de medição para evitar que queda de tensão em trilhas seja interpretada como sinal.

\subsection{Desacoplamento e transientes}
Na folha de alimentação, C1 de 47\,$\mu$F aparece entre 5 V e GND na saída do buck. Na folha de periféricos, há 100\,$\mu$F e 100\,nF junto da alimentação de 3,3 V do módulo microSD. O capacitor maior fornece reserva local para variações mais lentas de corrente; o cerâmico de 100\,nF fornece um caminho de baixa impedância para componentes rápidos. Eles se complementam: um valor grande isolado não substitui um capacitor pequeno colocado perto do pino, e o capacitor pequeno não sustenta uma demanda longa.

Posicionar desacoplamentos com trilhas curtas de alimentação e retorno; medir afundamento e ruído nos pinos durante Wi-Fi, escrita SD e acionamentos. Verificar tensão nominal, corrente de ripple, polaridade e ESR (resistência série equivalente) conforme o conversor e o regulador efetivamente usados. Os módulos ESP e sensores podem incluir capacitores próprios; sua presença e valor não devem ser inferidos apenas pelo símbolo do conector. As diretrizes de hardware do ESP32-S3 recomendam desacoplamento local também no domínio de 3,3 V e junto à memória.

\subsection{Barramento \texorpdfstring{I\textsuperscript{2}C}{I2C} e resistores de pull-up}
I\textsuperscript{2}C usa SDA (dados) e SCL (clock) compartilhados. As saídas são do tipo \emph{dreno aberto}: um dispositivo consegue puxar a linha para zero, mas não a força ativamente para nível alto. Um resistor de \emph{pull-up} em cada linha fornece esse alto quando ninguém a puxa para baixo. No esquema v4, R5 e R6 são de 4,7\,k$\Omega$ para 3,3 V; o firmware configura 400\,kHz no mestre. MPU6050 e OLED compartilham o barramento; o conector J8 reserva SDA/SCL para o módulo associado ao INA228. Endereços diferentes permitem coexistência, mas o endereço efetivo do INA228 depende da montagem de A0/A1 no módulo e ainda deve ser registrado.

Módulos comerciais frequentemente já trazem resistores de pull-up. Todos os resistores da mesma linha ficam em paralelo:
\begin{equation}
\frac{1}{R_{\mathrm{eq}}}=\sum_{k=1}^{n}\frac{1}{R_k}.
\label{eq:pullup}
\end{equation}
Duas unidades de 4,7\,k$\Omega$ equivalem a 2,35\,k$\Omega$; quatro equivalem a 1,175\,k$\Omega$; cinco, a 940\,$\Omega$. Resistência excessivamente baixa exige que os dispositivos drenem corrente alta para produzir zero; resistência alta demais, junto à capacitância dos cabos e módulos, faz a borda subir lentamente. A janela admissível depende de $V_{\mathrm{DD}}$, corrente de afundamento, capacitância total e tempo de subida. O limite inferior aproximado é $R_{p,\min}=(V_{\mathrm{DD}}-V_{OL,\max})/I_{OL}$: se o dispositivo suportar 3\,mA com nível baixo máximo de 0,4\,V a 3,3\,V, resulta cerca de 967\,$\Omega$. Cinco resistores de 4,7\,k$\Omega$ em paralelo já ficariam abaixo desse exemplo. Para modo rápido a 400\,kHz, a nota de aplicação da TI dá $t_r \simeq 0{,}8473R_{\mathrm{eq}}C_b$ e limite de subida de 300\,ns. Se $C_b=100$\,pF, o limite superior calculado é cerca de 3,54\,k$\Omega$; 4,7\,k$\Omega$ sozinho pode ser lento, enquanto 2,35\,k$\Omega$ é um candidato. Em $C_b=200$\,pF, o limite superior cai para 1,77\,k$\Omega$. Esses exemplos não substituem medir o barramento real.

Antes de montar todos os módulos, identificar os pull-ups existentes em cada placa, medir a resistência equivalente de SDA e SCL com a alimentação desligada e decidir quais resistores de módulo devem ser removidos ou desconectados. \textbf{Manter ao menos um par de pull-ups a 3,3 V}, mas remover os pares excedentes quando a associação em paralelo sair da janela calculada. Se um módulo puxar SDA/SCL para 5 V, remover seus pull-ups para 5 V ou instalar adaptação de nível apropriada antes de ligar ao ESP32-S3. Confirmar o resultado no osciloscópio, com todos os módulos e cabos, durante transferências a 400\,kHz. Resistores dos encoders são outro circuito e não entram nessa conta.

\subsection{INA228 e shunt de 1 \texorpdfstring{m$\Omega$}{mΩ}}
O conector J8 da folha de periféricos é denominado INA228 e entradas do motorista; ele traz alimentação indicada como 5 V, GND, SDA/SCL e outros sinais. O firmware mestre contém um driver para o INA228 no endereço I\textsuperscript{2}C \texttt{0x40} e publica grandezas elétricas quando a leitura é válida. O código configura shunt nominal de 1\,m$\Omega$ e corrente máxima de 75\,A. A posição física do shunt, o módulo específico e um ensaio de calibração ainda precisam ser confirmados na montagem. Os parâmetros a seguir descrevem a intenção de projeto e a configuração do software, não uma leitura experimental já comprovada.

\begin{figure}[htbp]\centering
\begin{tikzpicture}[>=Latex,bloco/.style={draw=nexus,rounded corners,fill=pale,align=center,minimum height=9mm,font=\small}]
\node[bloco,text width=2.6cm] (fonte) at (0,0) {Banco de baterias\\positivo};
\node[bloco,text width=2.0cm] (shunt) at (3.5,0) {Shunt\\1 m$\Omega$};
\node[bloco,text width=2.5cm] (carga) at (7.0,0) {Controlador\\do motor};
\node[bloco,text width=2.4cm] (ina) at (3.5,-2.2) {INA228\\interface I\textsuperscript{2}C};
\draw[->,very thick] (fonte)--(shunt);
\draw[->,very thick] (shunt)--(carga);
\draw[->,accent,thick] (shunt.south west)--node[left,font=\scriptsize]{IN+}(ina.north west);
\draw[->,accent,thick] (shunt.south east)--node[right,font=\scriptsize]{IN$-$}(ina.north east);
\node[bloco,text width=2.5cm] (vbus) at (7.0,-2.2) {Ponto VBUS\\após o shunt};
\draw[->,accent,dashed] (carga.south)--(vbus.north);
\draw[->,accent,dashed] (vbus.west)--(ina.east);
\end{tikzpicture}
\caption{Exemplo de medição no lado positivo com VBUS após o shunt. As tomadas Kelvin e o ponto VBUS devem ser definidos no projeto elétrico; o desenho não comprova o chicote instalado.}\label{fig:shunt-kelvin}
\end{figure}

O INA228 recebe duas medições distintas. \textbf{VSHUNT} é a diferença entre IN+ e IN$-$, tomada diretamente nos extremos do elemento resistivo por fios Kelvin; dividida pela resistência conhecida, fornece a corrente e seu sentido. \textbf{VBUS} é a tensão entre o pino VBUS e GND no ponto positivo escolhido do barramento do banco de baterias. Não é a queda no shunt: ela indica a tensão do barramento que alimenta o controlador. Se o pino for conectado antes ou depois do shunt, a leitura difere aproximadamente pela própria queda no shunt; a posição precisa constar do esquema final. A potência elétrica pode ser calculada a partir de corrente e tensão de barramento coerentes no tempo, $P\simeq V_{\mathrm{BUS}}I$, com convenção de sinal definida para consumo e regeneração. O driver no Mestre já lê os registradores de VSHUNT, VBUS, corrente, potência, temperatura do chip, energia e carga acumuladas. A validade exige resposta do dispositivo e verificações de faixa; uma leitura disponível ainda requer calibração experimental da montagem.

Com $R_{\mathrm{shunt}}=1\,\mathrm{m}\Omega=0{,}001\,\Omega$:
\begin{equation}
V_{\mathrm{shunt}}=IR_{\mathrm{shunt}},\qquad
1\,\mathrm{A}\longleftrightarrow 1\,\mathrm{mV},\qquad
75\,\mathrm{A}\longleftrightarrow 75\,\mathrm{mV}.
\label{eq:shunt}
\end{equation}
O calor no elemento resistivo é $P_{\mathrm{shunt}}=I^2R$: em 75 A, 5,625 W. Um shunt anunciado para 75 A precisa suportar essa dissipação no ambiente e na fixação reais, com margem térmica, seção dos terminais e tempo de sobrecarga adequados. A resistência muda com temperatura e com a tolerância de fabricação. As duas linhas de medição devem chegar a pontos Kelvin próximos ao elemento resistivo, separados do trajeto de corrente de potência, para não somar queda de tensão dos parafusos e cabos.

Segundo a folha de dados da TI, o INA228 oferece faixas diferenciais de $\pm163{,}84$\,mV e $\pm40{,}96$\,mV em VSHUNT. Para 75 A com shunt de 1\,m$\Omega$, a faixa de $\pm40{,}96$\,mV saturaria antes da corrente nominal; deve-se selecionar a faixa ampla (\texttt{ADCRANGE=0}) se a meta for medir até 75 A. Esse nome identifica a seleção do conversor interno do INA228. VBUS aceita de 0 a 85 V segundo a folha de dados do circuito integrado; a tensão máxima do banco de baterias, inclusive sob carga completa e transitórios, precisa respeitar o limite do circuito e do módulo comprado. A alimentação lógica VS e os pull-ups SDA/SCL são assuntos distintos: o conector J8 indica 5 V, mas isso não autoriza aplicar 5 V aos GPIOs do mestre. Verificar regulador, tradutor de nível e resistores da placa INA228 adquirida.

Para converter o registrador CURRENT em ampères, o firmware configura a calibração do INA228 com a resistência nominal do shunt. A folha de dados define, para \texttt{ADCRANGE=0},
\begin{equation}
\mathrm{CURRENT\_LSB}\geq \frac{I_{\max}}{2^{19}},
\qquad
\mathrm{SHUNT\_CAL}=13107{,}2\cdot10^6\,
       \mathrm{CURRENT\_LSB}\,R_{\mathrm{shunt}} .
\label{eq:ina-cal}
\end{equation}
No código, \texttt{CURRENT\_LSB} é calculado como $75/2^{19}$\,A por bit, aproximadamente 143\,$\mu$A, e \texttt{SHUNT\_CAL} resulta em cerca de 1875 para 1\,m$\Omega$. O Mestre confere identificação do chip e leitura de retorno da calibração, usa a faixa ampla de VSHUNT e rejeita erro matemático, conversão ainda não pronta, temperatura fora da faixa e corrente acima do limite admitido. Após três erros seguidos, suspende a leitura e tenta reinicializar o módulo após cinco segundos. A telemetria marca grandezas indisponíveis quando não há leitura válida; o simulador não deve ser confundido com o sensor. Definir a polaridade IN+/IN$-$ para que consumo e regeneração tenham sinais inequívocos. Comparar leituras em vários pontos com instrumento de referência, inclusive perto de zero e perto do máximo. Um alerta de sobrecorrente por software não substitui proteção elétrica independente.

\subsection{Relé habilitador, botão de emergência e diodo de roda livre}
A função prevista para K1 é interromper a \emph{linha de habilitação} do driver do motor elétrico, e não conduzir a corrente de tração. A ligação prevista usa o contato normalmente aberto (NA): a bobina precisa estar energizada para habilitar o driver; sem energia, o contato abre. Um botão físico de emergência também precisa interrompê-la independentemente dos ESPs. O diagrama abaixo representa essa intenção, não o circuito da PCB nem um chicote validado. Um resistor de pull-down externo na base de Q4 deve manter a bobina desligada durante reset e antes da configuração do GPIO. O nível elétrico exigido pela entrada de habilitação e o comportamento real do driver precisam ser medidos e confrontados com sua documentação.

\begin{figure}[htbp]\centering
\begin{tikzpicture}[>={Latex[length=2mm]},
  bloco/.style={draw=nexus,rounded corners=2mm,fill=pale,align=center,minimum height=11mm,font=\small,inner sep=5pt},
  habilitacao/.style={->,line width=0.9pt,draw=nexus},
  acionamento/.style={->,line width=0.9pt,draw=accent}]
\node[bloco,text width=2.3cm] (sinal) at (0,0) {Fonte do sinal\\de habilitação};
\node[bloco,text width=2.5cm] (botao) at (3.4,0) {Botão de emergência\\NF};
\node[bloco,text width=2.3cm] (contato) at (6.8,0) {Contato K1\\NA previsto};
\node[bloco,text width=2.4cm] (driver) at (10.2,0) {Entrada ENABLE\\do driver};
\draw[habilitacao] (sinal)--(botao);
\draw[habilitacao] (botao)--(contato);
\draw[habilitacao] (contato)--(driver);
\node[bloco,text width=2.3cm] (gpio) at (0,-2.0) {Saída do ESP\\futura};
\node[bloco,text width=2.5cm] (base) at (3.4,-2.0) {Resistor de base\\e estado seguro};
\node[bloco,text width=2.3cm] (trans) at (6.8,-2.0) {Transistor\\Q4};
\node[bloco,text width=2.4cm] (bobina) at (10.2,-2.0) {Bobina K1\\D1 em paralelo};
\draw[acionamento] (gpio)--(base);
\draw[acionamento] (base)--(trans);
\draw[acionamento] (trans)--(bobina);
\draw[acionamento,densely dashed] (bobina.north)--++(0,0.52) -| (contato.south);
\end{tikzpicture}
\caption{Integração prevista: botão de emergência NF e contato K1 NA em série na linha de habilitação. A bobina energizada fecha o contato; sem energia o driver fica desabilitado. A seta tracejada relaciona bobina e contato.}\label{fig:rele-habilitador}
\end{figure}

A folha do esquema contém K1 (símbolo SANYOU SRD), Q4 (BC337), R1 de 1\,k$\Omega$, R2 de 10\,k$\Omega$ e D1 (1N4007). O Mestre comanda Q4 pelo GPIO21. O armamento energiza K1 e a parada desenergiza a bobina e comanda 0 V no DAC. Confira o estado do contato na placa e o tempo de resposta durante falha e reinício. O botão físico deve funcionar independentemente do navegador, Wi-Fi e UART.

A opção pelo contato normalmente aberto coincide com o arquivo \path{hardware/kicad/KiCad_pcb_telemetria/KiCad_pcb_telemetria.kicad_pcb}: nele, o comum de K1 recebe \texttt{+48V}, o contato normalmente aberto alimenta \texttt{DRV\_ENABLE\_SW} e o normalmente fechado fica sem conexão. A PCB consultada também liga \texttt{DRV\_MODE\_B\_CMD}, \texttt{DRV\_MODE\_A\_CMD} e \texttt{DRV\_ENABLE\_CMD} aos GPIOs 13, 14 e 21, usados para atuação pelo firmware.

A bobina é indutiva. Ao interromper sua corrente, o campo magnético tenta mantê-la e pode produzir sobretensão no transistor. O diodo de roda livre D1 em paralelo oferece um caminho temporário para essa corrente: cátodo no lado positivo da bobina e ânodo no lado comutado pelo transistor, ficando reversamente polarizado durante o acionamento normal. Ele reduz o pico de tensão, mas pode aumentar o tempo de soltura do relé; medir esse tempo no circuito final. Q4 evita exigir do GPIO a corrente da bobina; R1 limita a corrente de base, e o estado de desligamento da base precisa ser definido durante reset e partida.

\begin{nota}{Conferência obrigatória da folha v4}
No desenho exportado da folha, R2 de 10\,k$\Omega$ aparece em série com a alimentação da bobina e Q4 como derivação desse ramo; essa ligação não corresponde diretamente ao acionamento por transistor em série mostrado funcionalmente acima. A queda de tensão em R2 pode impedir a corrente nominal da bobina. É necessário revisar a ligação no esquema, conferir bobina, contato escolhido, orientação de D1 e estado sem alimentação, executar a checagem elétrica do KiCad e testar com fonte limitada em bancada antes de conectar o driver. A posição do botão de emergência em série também deve constar do esquema e do chicote finais.
\end{nota}

Outros transistores da folha representam sinais de modo e reversão e o MCP4725 aparece como DAC (conversor digital-analógico) para a referência de acelerador. O MCP4725 inicia em 0 V e só aplica a referência de tração quando o controle habilita o relé. Antes de integração, mapear tensão exigida pelo driver, estado seguro de cada fio, isolação, tempo de resposta, limites físicos e retorno medido; testar separadamente de toda a tração.

\paragraph{Fontes técnicas.}
As faixas e equações do INA228 vêm da \href{https://www.ti.com/lit/ds/symlink/ina228.pdf}{folha de dados da Texas Instruments}. O cálculo da janela de pull-ups segue \href{https://www.ti.com/lit/an/slva689/slva689.pdf}{I\textsuperscript{2}C Bus Pull-Up Resistor Calculation}, da TI. O desacoplamento do ESP deve ser confrontado com o \href{https://docs.espressif.com/projects/esp-hardware-design-guidelines/en/latest/esp32s3/schematic-checklist.html}{checklist de esquema do ESP32-S3}, da Espressif, e com a folha de dados de cada módulo comprado.
